% Minimal preamble for Open Logic Project sources (pinned commit 1e960beff9ed7835bf3e3f1335e21af3439cd107).
%
% OLP's own master preamble cannot run on the project host: TeX Live 2025 there lacks
% prooftree.sty, ntheorem.sty, extarrows.sty, mathdots.sty, isomath.sty and subfiles.sty.
% Following the project's own rule ("if the author preamble needs packages the toolchain lacks,
% do not invent macros -- use a minimal working preamble"), this file declares, verbatim in
% spirit and numbering behaviour, exactly the theorem-like environments OLP defines in
% sty/open-logic-book-envs.sty (thm/lem/prop/cor/defn/ex/prob/rem/note/case/axiom/conv),
% plus the shorthand accents OLP uses (`!!a`, `!!b`, `!!{...}`).
%
% Author text is untouched; only the packaging is ours.
\documentclass[10pt,a4paper]{article}

\usepackage{amsmath}
\usepackage{amssymb}
\usepackage{amsthm}
\usepackage{graphicx}
\usepackage{hyperref}
\usepackage{enumitem}
\usepackage[framemethod=default]{mdframed}

% ---- OLP shorthand accents (from sty/open-logic.sty) ----
% NB: \! is taken by LaTeX (negative thin space), so we only define the !!a / !!b families
% and the \olup helper used by them.
\newcommand*{\olup}[1]{\MakeUppercase{#1}}
\def\!!!a#1{\olup{#1}}
\def\!!a#1{\olup{#1}}
\def\!!b#1{\olup{#1}}

% ---- theorem-like environments, same family numbering as OLP ----
\theoremstyle{plain}
\newtheorem{thm}{Theorem}[section]
\newtheorem{lem}[thm]{Lemma}
\newtheorem{prop}[thm]{Proposition}
\newtheorem{cor}[thm]{Corollary}
\newtheorem{axiom}{Axiom}

\theoremstyle{definition}
\newtheorem{defn}[thm]{Definition}
\newtheorem{ex}[thm]{Example}
\newtheorem{prob}{Problem}[section]
\newtheorem{probd}{Problem}
\newtheorem{conv}[thm]{Convention}

\theoremstyle{remark}
\newtheorem{rem}[thm]{Remark}
\newtheorem{note}[thm]{Note}
\newtheorem{case}[thm]{Case}

% ---- environments OLP defines as wrappers ----
\newenvironment{probtag}[1]{\begin{prob}}{\end{prob}}
\newenvironment{explain}{}{}
\newenvironment{explain*}{}{}
\newenvironment{editorial}{}{}
\newenvironment{digress}{}{}
\newenvironment{digress*}{}{}
\newenvironment{history}{}{}
\newenvironment{reference}{}{}
\newcommand*{\olfileid}[3]{}
\newcommand*{\olsection}[1]{\section*{#1}}
\newcommand*{\olchapter}[3]{}
\newcommand*{\olpart}[2]{}
\newcommand*{\olimport}[1]{}
\newcommand*{\olinclude}[1]{}
\newcommand*{\olpath}[1]{}
\newcommand*{\ollangid}[1]{}
\newcommand*{\problemtag}[1]{}
\newcommand*{\problemauthor}[1]{}
\newcommand*{\ollabel}[1]{\label{#1}}
% \olref takes 0-3 OPTIONAL arguments plus one mandatory label:
%     \olref{lab}  |  \olref[from]{lab}  |  \olref[from][to]{lab}  |  \olref[a][b][c]{lab}
% A single-argument definition silently swallows "[sth]" and emits a broken reference
% (measured on the host: the text came out as "??iso]x"). The chain below consumes the
% brackets and cross-references only the mandatory argument.
\makeatletter
\newcommand{\olref}{\@ifnextchar[{\olref@i}{\olref@final}}
\def\olref@i[#1]{\@ifnextchar[{\olref@ii}{\olref@final}}
\def\olref@ii[#1]{\@ifnextchar[{\olref@iii}{\olref@final}}
\def\olref@iii[#1]{\olref@final}
\def\olref@final#1{\ref{#1}}
\makeatother
\newcommand*{\olnewpage}{\clearpage}

% ===================== OLP compatibility layer (1) =====================
% Compatibility layer for Open Logic Project sources (pinned commit 1e960beff9ed7835bf3e3f1335e21af3439cd107).
%
% OLP's own style stack (sty/open-logic.sty) needs ntheorem, prooftree, extarrows, mathdots,
% isomath and subfiles; the project's local and host TeX Live installs lack several of those.
% Rather than vendor packages, this file supplies drop-in equivalents for every OLP construct
% that the frozen units actually use. It only rearranges the author's own tokens; no
% mathematical content is added, and nothing here is claimed to be OLP's own file.

\makeatletter

% ---------------- proof trees (OLP's prooftree is a bussproofs derivative) ----------------
% These commands CONSUME their argument and typeset nothing: the bracket labels of OLP's proof
% trees carry rule names. Printing them inline dropped those labels into math mode and produced
% 63 "Missing $" / "allowed only in math mode" errors; swallowing them keeps the brace balance
% and the mode correct. The derivation text between the commands is still printed unchanged.
\providecommand{\AxiomC}[1]{}
\providecommand{\UnaryInfC}[1]{}
\providecommand{\BinaryInfC}[1]{}
\providecommand{\TrinaryInfC}[1]{}
\providecommand{\QuaternaryInfC}[1]{}
\providecommand{\QuinaryInfC}[1]{}
\providecommand{\RightLabel}[1]{}
\providecommand{\LeftLabel}[1]{}
\providecommand{\doubleLine}{}
\providecommand{\noLine}{}
\providecommand{\DisplayProof}{}
\providecommand{\fCenter}{}
\newenvironment{prooftree}{}{}
\newenvironment{ltprooftree}{}{}
\newenvironment{oltableau}{}{}

% ---------------- logical notation ----------------
\providecommand{\lif}{\mathbin{\rightarrow}}
\providecommand{\lforall}{\mathord{\forall}}
\providecommand{\lexists}{\mathord{\exists}}
\providecommand{\ltrue}{\top}
\providecommand{\lfalse}{\bot}
\providecommand{\Ltrue}{\top}
\providecommand{\Lfalse}{\bot}
\providecommand{\Sequent}{\mathrel{\Rightarrow}}
\providecommand{\Proves}{\mathrel{\vdash}}
\providecommand{\Entails}{\mathrel{\models}}
\providecommand{\Subst}[3]{#1[#2/#3]}
\providecommand{\sFmla}[2]{#1(#2)}
\providecommand{\pFmla}[2]{#1(#2)}
\providecommand{\formula}{\relax}
\providecommand{\True}{\top}
\providecommand{\False}{\bot}
\providecommand{\Args}[1]{#1}
\providecommand{\Ts}{\ensuremath{\top}}
\providecommand{\Fi}{\ensuremath{\bot}}
\providecommand{\lfalsum}{\bot}

% ---------------- objects, structures, metatheory ----------------
\providecommand{\Obj}{}
\providecommand{\Struct}{\ensuremath{\mathfrak{A}}}
\providecommand{\Domain}{\ensuremath{|\mathfrak{A}|}}
\providecommand{\Assign}{\ensuremath{s}}
\providecommand{\Sat}[2]{\ensuremath{#1 \models #2}}
\providecommand{\mSat}[3]{\ensuremath{#1 \models #2[#3]}}
\providecommand{\mModel}[2]{\ensuremath{#1 \models #2}}
\providecommand{\Lang}{\ensuremath{\mathcal{L}}}
\providecommand{\Th}{\ensuremath{\mathrm{Th}}}
\providecommand{\Log}{\ensuremath{\mathrm{L}}}
\providecommand{\tuple}[1]{\langle #1 \rangle}
\providecommand{\num}[1]{\underline{#1}}
\providecommand{\cfind}[1]{\varphi_{#1}}
\providecommand{\Nat}{\ensuremath{\mathbb{N}}}
\providecommand{\Mult}{\ensuremath{\mathrm{Mult}}}
\providecommand{\Add}{\ensuremath{\mathrm{Add}}}
\providecommand{\Exp}{\ensuremath{\mathrm{Exp}}}
\providecommand{\Zero}{\ensuremath{0}}
\providecommand{\Succ}{\ensuremath{\mathrm{S}}}
\providecommand{\lambd}{\ensuremath{\lambda}}
\providecommand{\bredpar}{\mathrel{\rightarrow_\beta}}
\providecommand{\bcd}[1]{\ensuremath{\ulcorner #1 \urcorner}}
\providecommand{\Set}[1]{\ensuremath{\{#1\}}}
\providecommand{\indcase}[1]{\par\noindent\emph{Case #1}.\quad}
\providecommand{\indfrm}[1]{\par\noindent\emph{Form #1}.\quad}
\providecommand{\iftag}[1]{#1}
\providecommand{\tagitem}[1]{\item[#1]}
% these five are used with arguments by OLP's proof trees -> swallow the argument
\providecommand{\Elim}[1]{}
\providecommand{\Intro}[1]{}
\providecommand{\LeftR}[1]{}
\providecommand{\RightR}[1]{}
\providecommand{\TAss}[1]{}
\providecommand{\TRule}[1]{}

\makeatother
\providecommand{\olpath}{.}

% ===================== OLP compatibility layer (2) =====================
% --- second compatibility block: notation and shorthands used by the frozen OLP units ---
% Everything here is \providecommand, so anything already defined by LaTeX or the packages wins.
\providecommand{\CutCS}{\mathrm{Cut}}
\providecommand{\LeftSubproofLabel}{\relax}
\providecommand{\RightSubproofLabel}{\relax}
\providecommand{\maeh}{\relax}
\providecommand{\cheight}{\relax}
\providecommand{\ande}{\mathrm{and}_E}
\providecommand{\andi}{\mathrm{and}_I}
\providecommand{\ore}{\mathrm{or}_E}
\providecommand{\ori}{\mathrm{or}_I}
\providecommand{\typeof}{\mathord{:}}
\providecommand{\pValue}[1]{#1}
\providecommand{\Sent}{\mathrm{Sent}}
\providecommand{\cardplus}{\mathbin{+}}
\providecommand{\cardtimes}{\mathbin{\cdot}}
\providecommand{\cardexpo}{\mathbin{\hat{\ }}}
\providecommand{\cardnless}{\not<}
\providecommand{\cardle}{\leq}
\providecommand{\funfromto}[3]{#1\colon #2 \to #3}
\providecommand{\funimage}[1]{\operatorname{im}(#1)}
\providecommand{\fdefined}{\mathord{\downarrow}}
\providecommand{\canonord}{\operatorname{ord}}
\providecommand{\ordtimes}{\mathbin{\cdot}}
\providecommand{\ordeq}[2]{#1 = #2}
\providecommand{\ordneq}[2]{#1 \neq #2}
\providecommand{\ordtype}[1]{\operatorname{ord}(#1)}
\providecommand{\ordsucc}{\mathord{+1}}
\providecommand{\ordunion}{\cup}
\providecommand{\onesphere}{S^1}
\providecommand{\rotationsgroup}{\mathrm{SO}(2)}
\providecommand{\equivrep}[1]{[#1]}
\providecommand{\rlexless}{\prec}
\providecommand{\ZFC}{\mathrm{ZFC}}
\providecommand{\ZF}{\mathrm{ZF}}
\providecommand{\Rat}{\mathbb{Q}}
\providecommand{\Real}{\mathbb{R}}
\providecommand{\PosInt}{\mathbb{Z}^{+}}
\providecommand{\Id}{\mathrm{id}}
\providecommand{\Bin}{\mathrm{Bin}}
\providecommand{\pto}{\rightharpoonup}
\providecommand{\dom}{\operatorname{dom}}
\providecommand{\ran}{\operatorname{ran}}
\providecommand{\len}{\operatorname{len}}
\providecommand{\concat}{\mathbin{{}^\frown}}
\providecommand{\card}[1]{\lvert #1 \rvert}
\providecommand{\Pow}{\mathcal{P}}
\providecommand{\Setabs}[2]{\{\, #1 \mid #2 \,\}}
\providecommand{\comp}{\circ}
\providecommand{\Char}{\chi}
\providecommand{\Prf}{\operatorname{Prf}}
\providecommand{\Prov}{\operatorname{Prov}}
\providecommand{\Proj}{\operatorname{proj}}
\providecommand{\Gn}{\operatorname{gn}}
\providecommand{\gn}{\operatorname{gn}}
\providecommand{\Value}[1]{\operatorname{val}(#1)}
\providecommand{\Deduce}{\mathrm{D}}
\providecommand{\Dual}{\mathrm{d}}
\providecommand{\Ax}{\mathrm{Ax}}
\providecommand{\pSat}[3]{#1 \models #2\,[#3]}
\providecommand{\pAssign}{\sigma}
\providecommand{\bforall}{\forall}
\providecommand{\bexists}{\exists}
\providecommand{\liff}{\mathbin{\leftrightarrow}}
\providecommand{\FV}{\operatorname{FV}}
\providecommand{\UnaryInf}{\ensuremath{\mid}}
\providecommand{\BinaryInf}{\ensuremath{\mid}}
\providecommand{\depth}{\operatorname{depth}}
\providecommand{\oliflabeldef}[3]{#1}
\providecommand{\TMblank}{\sqcup}
\providecommand{\TMstroke}{\mid}
\providecommand{\TMstay}{\mathrm{S}}
\providecommand{\TMright}{\mathrm{R}}
\providecommand{\TMleft}{\mathrm{L}}
\providecommand{\TMendtape}{\triangleright}
\providecommand{\umin}{\mu}
\providecommand{\Sum}{\sum}
% bibliography helpers (natbib): the frozen units cite the shared OLP bibliography
\providecommand{\citep}[1]{\cite{#1}}
\providecommand{\citet}[1]{\cite{#1}}
\providecommand{\citeauthor}[1]{\cite{#1}}
\providecommand{\citeyear}[1]{\cite{#1}}

% ===================== OLP compatibility layer (3) =====================
% --- third compatibility block: remaining notation seen in the frozen OLP units ---
\providecommand{\iso}{\cong}
\providecommand{\nszero}{\mathrm{0}}
\providecommand{\nssucc}{\mathrm{S}}
\providecommand{\nsplus}{\mathbin{+}}
\providecommand{\nssless}{<}
\providecommand{\nsless}{<}
\providecommand{\PL}{\mathrm{PL}}
\providecommand{\Nec}{\mathord{\Box}}
\providecommand{\RK}{\mathrm{RK}}
\providecommand{\ST}{\mathrm{ST}}
\providecommand{\SSubst}[3]{#1[#2/#3]}
\providecommand{\subst}[3]{#1[#2/#3]}
\providecommand{\tagfalse}{\mathrm{F}}
\providecommand{\cutr}{\mathrm{Cut}}
\providecommand{\cardneq}{\neq}
\providecommand{\issucc}[1]{\operatorname{succ}(#1)}
\providecommand{\tup}[1]{\langle #1 \rangle}
\providecommand{\Var}{\mathrm{Var}}
\providecommand{\Bew}{\operatorname{Bew}}
\providecommand{\Con}{\operatorname{Con}}
\providecommand{\Cons}{\operatorname{Cons}}
\providecommand{\pDeduce}{\mathrm{D}}
\providecommand{\psat}{\models}
\providecommand{\den}[1]{\llbracket #1 \rrbracket}
\providecommand{\Sem}[2]{\llbracket #1 \rrbracket_{#2}}
\providecommand{\DomainOf}[1]{\lvert #1 \rvert}
\providecommand{\Form}{\mathrm{Form}}
\providecommand{\Terms}{\mathrm{Term}}
\providecommand{\Atom}{\mathrm{Atom}}
\providecommand{\FVars}[1]{\operatorname{FV}(#1)}
\providecommand{\fv}[1]{\operatorname{fv}(#1)}
\providecommand{\tuplep}[1]{\langle #1 \rangle}
\providecommand{\restrict}[2]{#1 \upharpoonright #2}
\providecommand{\sto}{\rightarrow}
\providecommand{\FmlaSet}[1]{\mathrm{Fmla}_{#1}}
\providecommand{\Thm}{\mathrm{Thm}}
\providecommand{\Eval}{\operatorname{val}}
\providecommand{\blank}{\sqcup}
\providecommand{\blankTape}{\sqcup}
\providecommand{\Consis}{\operatorname{Con}}
\providecommand{\drs}{\mathrm{d}}
\providecommand{\VRule}{\mid}
\providecommand{\Bet}{\mathrm{B}}
\providecommand{\Seg}[1]{\mathrm{Seg}_{#1}}
\providecommand{\Enums}{\mathrm{En}}
\providecommand{\Z}{\mathbb{Z}}
\providecommand{\N}{\mathbb{N}}
\providecommand{\abs}[1]{\lvert #1 \rvert}
\providecommand{\floor}[1]{\lfloor #1 \rfloor}
\providecommand{\ceil}[1]{\lceil #1 \rceil}

% ===================== OLP compatibility layer (4) =====================
% --- fourth compatibility block: harvested from the per-candidate compile logs ---
\providecommand{\beredpar}{\mathrel{\rightarrow_\beta}}
% \bered / \bred are used as infix relations AND with a bracketed annotation, so they need the
% optional-argument chain; \providecommand would leave an argument-less definition in place.
\makeatletter
\def\bered{\@ifnextchar[{\bered@i}{\mathrel{\rightarrow_\beta}}}
\def\bered@i[#1]{\mathrel{\rightarrow_\beta}}
\def\bred{\@ifnextchar[{\bred@i}{\mathrel{\rightarrow_\beta}}}
\def\bred@i[#1]{\mathrel{\rightarrow_\beta}}
\makeatother
\providecommand{\beredone}{\mathrel{=_\beta}}
\providecommand{\bredone}{\mathrel{=_\beta}}
\providecommand{\xredone}{\mathrel{=_\beta}}
\providecommand{\aconvone}{\mathrel{=_\alpha}}
\providecommand{\cardless}{<}
\providecommand{\tagrefs}[1]{#1}
\providecommand{\ordplus}{\mathbin{+}}
\providecommand{\funrestrictionto}[2]{#1 \upharpoonright #2}
\providecommand{\fundefined}{\mathord{\uparrow}}
\providecommand{\Trm}{\mathrm{Trm}}
\providecommand{\OProv}{\operatorname{Prov}}
\providecommand{\OPrf}{\operatorname{Prf}}
\providecommand{\aeq}{\mathrel{=}}
\providecommand{\mClass}[1]{[#1]}
\providecommand{\equivclass}[1]{[#1]}
\providecommand{\pheight}{\operatorname{ht}}
\providecommand{\phy}{\operatorname{ht}}
\providecommand{\Cut}{\mathrm{Cut}}
\providecommand{\QuantRank}{\operatorname{qr}}
\providecommand{\elemequiv}{\mathrel{\equiv}}
\providecommand{\PIso}{\mathrm{PIso}}
\providecommand{\Theory}{\mathrm{Th}}
\providecommand{\supstrict}{\subsetneq}
\providecommand{\trcl}{\operatorname{trcl}}
\providecommand{\cref}[1]{\ref{#1}}
\providecommand{\Discharge}{\mathrm{Dis}}
\providecommand{\Weakening}{\mathrm{Wk}}
\providecommand{\ZFminus}{\mathrm{ZF}^{-}}
\providecommand{\RK}{\mathrm{RK}}
% \Dual is printed both in math mode and (with a bracketed annotation) in running text;
% the math-mode test keeps \mathrm legal in both cases.
\makeatletter
\renewcommand{\Dual}{\@ifnextchar[{\Dual@i}{\ifmmode\mathrm{d}\else$\mathrm{d}$\fi}}
\def\Dual@i[#1]{\ifmmode\mathrm{d}[#1]\else$\mathrm{d}[#1]$\fi}
\makeatother
\providecommand{\RightR}{\mathrm{R}}
\providecommand{\LeftR}{\mathrm{L}}
\providecommand{\Intro}{\mathrm{I}}
\providecommand{\Axiom}{\mathrm{Ax}}
\providecommand{\formula}[1]{#1}
\providecommand{\stared}{\ast}
\providecommand{\oset}[2]{\{\, #1 \mid #2 \,\}}
\providecommand{\setof}[2]{\{\, #1 \mid #2 \,\}}
\providecommand{\pow}[1]{\mathcal{P}(#1)}
\providecommand{\ordpair}[2]{\langle #1, #2 \rangle}
\providecommand{\ordtriple}[3]{\langle #1, #2, #3 \rangle}
\providecommand{\union}{\cup}
\providecommand{\inter}{\cap}
\providecommand{\seq}[1]{\langle #1 \rangle}
\providecommand{\SuchThat}{\mid}
\providecommand{\where}{\mid}
\providecommand{\TVar}{\mathrm{TVar}}
\providecommand{\Bool}{\mathrm{Bool}}
\providecommand{\Val}{\operatorname{Val}}
\providecommand{\FV}{\operatorname{FV}}
\providecommand{\TL}[1]{\mathrm{TL}}
\providecommand{\xstared}{\ast}

% ===================== OLP compatibility layer (5) =====================
% --- fifth compatibility block: rescues the quarantined OLP candidates ---
% The proof-tree and mode fixes live in block 1 (where those macros are first defined); this block
% only adds notation that the quarantined spans still needed and that no earlier block covered.
\providecommand{\ident}{\operatorname{id}}
\providecommand{\setrank}{\operatorname{rank}}
\providecommand{\cardfont}{\mathfrak{c}}
\providecommand{\Complement}{\overline}
\providecommand{\QR}{\mathrm{QR}}
\providecommand{\becd}{\ensuremath{\ulcorner}}
\providecommand{\disjointsum}{\mathbin{\oplus}}
\providecommand{\cardeq}[2]{#1 = #2}
\providecommand{\xred}{\mathrel{\rightarrow}}
\providecommand{\cardless}{<}
\providecommand{\cardplus}{\mathbin{+}}
\providecommand{\cardtimes}{\mathbin{\cdot}}
\providecommand{\cardexpo}{\mathbin{\hat{\ }}}
\providecommand{\cardneq}{\neq}
\providecommand{\cardle}{\leq}
\providecommand{\cardnless}{\not<}
\begin{document}

% ==== OLP file: content/normal-modal-logic/syntax-and-semantics/entailment.tex ====
% Part: normal-modal-logic
% Chapter: syntax-and-semantics
% Section: entailment



\olfileid{nml}{syn}{ent}

\olsection{Entailment}

\begin{explain}
  With the definition of truth at a world, we can define an entailment
  relation between !!{formula}s. !!^a{formula}~$!B$ entails~$!A$ iff,
  whenever $!B$ is true, $!A$ is true as well. Here, ``whenever''
  means both ``whichever model we consider'' as well as ``whichever
  world in that model we consider.''
\end{explain}

\begin{defn}
  If $\Gamma$ is a set of !!{formula}s and $!A$ !!a{formula}, then
  $\Gamma$ \emph{entails}~$!A$, in symbols: $\Gamma \Entails !A$, if
  and only if for every model $\mModel{M} = \tuple{W, R, V}$ and world
  $w \in W$, if $\mSat{M}{!B}[w]$ for every $!B \in \Gamma$, then
  $\mSat{M}{!A}[w]$. If $\Gamma$ contains a single !!{formula}~$!B$,
  then we write $!B \Entails !A$.
\end{defn}

\begin{ex}
  To show that !!a{formula} entails another, we have to reason about
  all models, using the definition of $\mSat{M}{}[w]$. For instance,
  to show $p \lif \Diamond p \Entails \Box\lnot p \lif \lnot p$, we
  might argue as follows: Consider a model $\mModel{M} = \tuple{W, R,
    V}$ and $w \in W$, and suppose $\mSat{M}{p \lif \Diamond
    p}[w]$. We have to show that $\mSat{M}{\Box\lnot p \lif \lnot
    p}[w]$. Suppose not. Then $\mSat{M}{\Box\lnot p}[w]$ and
  $\mSat/{M}{\lnot p}[w]$. Since $\mSat/{M}{\lnot p}[w]$, $\mSat{M}{
    p}[w]$. By assumption, $\mSat{M}{p \lif \Diamond p}[w]$, hence
  $\mSat{M}{\Diamond p}[w]$. By definition of $\mSat{M}{\Diamond
    p}[w]$, there is some $w'$ with $Rww'$ such that
  $\mSat{M}{p}[w']$. Since also $\mSat{M}{\Box \lnot p}[w]$,
  $\mSat{M}{\lnot p}[w']$, a contradiction.

  To show that !!a{formula}~$!B$ does not entail another~$!A$, we have
  to give a counterexample, i.e., a model $\mModel{M} = \tuple{W, R,
    V}$ where we show that at some world~$w \in W$, $\mSat{M}{!B}[w]$
  but $\mSat/{M}{!A}[w]$. Let's show that $p \lif \Diamond p \Entails/
  \Box p \lif p$. Consider the model in \olref{fig:counterex}.  We
  have $\mSat{M}{\Diamond p}[w_1]$ and hence $\mSat{M}{p \lif
    \Diamond p}[w_1]$. However, since $\mSat{M}{\Box p}[w_1]$ but
  $\mSat/{M}{p}[w_1]$, we have $\mSat/{M}{\Box p \lif p}[w_1]$.
  \begin{figure}
  \begin{center}
    \begin{tikzpicture}[modal]
      \node[world] (w1) [label=right:\mFalse{p}]{$w_1$}; 
      \node[world] (w2) [label=right:\mTrue{p}, above left=of w1]{$w_2$}; 
      \node[world] (w3) [label=right:\mTrue{p}, above right=of w1] {$w_3$};
      \draw[->] (w1) to (w2);
      \draw[->] (w1) to (w3);
    \end{tikzpicture}
  \end{center}
  \caption{Counterexample to $p \lif \Diamond p
  \Entails \Box p \lif p$.}
  \ollabel{fig:counterex}
  \end{figure}
  
  Often very simple counterexamples suffice. The model $\mModel{M'} =
  \{W', R', V'\}$ with $W' = \{w\}$, $R' = \emptyset$, and $V'(p) =
  \emptyset$ is also a counterexample: Since $\mSat/{M'}{p}[w]$,
  $\mSat{M'}{p \lif \Diamond p}[w]$. As no worlds are accessible
  from~$w$, we have $\mSat{M'}{\Box p}[w]$, and so $\mSat/{M'}{\Box p
    \lif p}[w]$.
\end{ex}

\begin{prob}
  Show that $\Box (!A \land !B) \Entails \Box !A$.
\end{prob}

\begin{prob}
  Show that $\Box (p \lif q) \Entails/ p \lif \Box q$ and $p \lif \Box
  q \Entails/\Box (p \lif q)$.
\end{prob}


% ==== OLP file: content/normal-modal-logic/syntax-and-semantics/introduction.tex ====
% Part: normal-modal-logic
% Chapter: syntax-and-semantics
% Section: introduction



\olfileid{nml}{syn}{int}

\olsection{Introduction}

Modal logic deals with \emph{modal propositions} and the entailment
relations among them. Examples of modal propositions are the
following:
\begin{enumerate}
\item It is necessary that $2+2=4$.
\item It is necessarily possible that it will rain tomorrow.
\item If it is necessarily possible that~$!A$ then it is possible
  that~$!A$.
\end{enumerate}
Possibility and necessity are not the only modalities: other unary
connectives are also classified as modalities, for instance, ``it
ought to be the case that~$!A$,'' ``It will be the case that~$!A$,''
``Dana knows that~$!A$,'' or ``Dana believes that~$!A$.''

Modal logic makes its first appearance in Aristotle's \emph{De
  Interpretatione}: he was the first to notice that necessity implies
possibility, but not vice versa; that possibility and necessity are
inter-definable; that If $!A \land !B$ is possibly true then
$!A$ is possibly true and $!B$ is possibly true, but not
conversely; and that if $!A \to !B$ is necessary, then if
$!A$ is necessary, so is~$!B$.

The first modern approach to modal logic was the work of C.~I. Lewis,
culminating with Lewis and Langford, \emph{Symbolic Logic}
(1932). Lewis \& Langford were unhappy with the representation of
implication by means of the material conditional: $!A \lif !B$
is a poor substitute for ``$!A$ implies $!B$.'' Instead, they
proposed to characterize implication as ``Necessarily, if $!A$
then $!B$,'' symbolized as $!A \strictif !B$. In trying to
sort out the different properties, Lewis identified five different
modal systems, \Log{S1}, \ldots, \Log{S4}, \Log{S5}, the last
two of which are still in use.

The approach of Lewis and Langford was purely \emph{syntactical}: they
identified reasonable axioms and rules and investigated what was
provable with those means. A semantic approach remained elusive for a
long time, until a first attempt was made by Rudolf Carnap in
\emph{Meaning and Necessity} (1947) using the notion of a \emph{state
  description}, i.e., a collection of atomic sentences (those that are
``true'' in that state description). After lifting the truth
definition to arbitrary sentences $!A$, Carnap defines $!A$
to be \emph{necessarily true} if it is true in all state
descriptions. Carnap's approach could not handle \emph{iterated}
modalities, in that sentences of the form ``Possibly necessarily
\ldots possibly $!A$'' always reduce to the innermost modality.

The major breakthrough in modal semantics came with Saul Kripke's
article ``A Completeness Theorem in Modal Logic'' (JSL 1959). Kripke
based his work on Leibniz's idea that a statement is necessarily true
if it is true ``at all possible worlds.'' This idea, though, suffers
from the same drawbacks as Carnap's, in that the truth of statement at
a world $w$ (or a state description $s$) does not depend on $w$ at
all. So Kripke assumed that worlds are related by an
\emph{accessibility relation} $R$, and that a statement of the form
``Necessarily $!A$'' is true at a world $w$ if and only if
$!A$ is true at all worlds $w'$ \emph{accessible from}
$w$. Semantics that provide some version of this approach are called
Kripke semantics and made possible the tumultuous development of modal
logics (in the plural).

When interpreted by the Kripke semantics, modal logic shows us what
\emph{relational structures} look like ``from the inside.'' A
relational structure is just a set equipped with a binary relation
(for instance, the set of students in the class ordered by their
social security number is a relational structure). But in fact
relational structures come in all sorts of domains: besides relative
possibility of states of the world, we can have epistemic states of
some agent related by epistemic possibility, or states of a dynamical
system with their state transitions, etc. Modal logic can be used to
model all of these: the first gives us ordinary, alethic, modal logic;
the others give us epistemic logic, dynamic logic, etc.

We focus on one particular angle, known to modal logicians as
``correspondence theory.'' One of the most significant early
discoveries of Kripke's is that many properties of the accessibility
relation~$R$ (whether it is transitive, symmetric, etc.)  can be
characterized \emph{in the modal language} itself by means of
appropriate ``modal schemas.'' Modal logicians say, for instance, that
the reflexivity of $R$ ``corresponds'' to the schema ``If
necessarily~$!A$, then~$!A$''. We explore mainly the correspondence
theory of a number of classical systems of modal logic (e.g., \Log{S4}
and \Log{S5}) obtained by a combination of the schemas \Ax{D}, \Ax{T},
\Ax{B}, \Ax{4}, and~\Ax{5}.


% ==== OLP file: content/normal-modal-logic/syntax-and-semantics/language-modal-logic.tex ====
% Part: normal-modal-logic
% Chapter: syntax-and-semantics
% Section: language-modal-logic



\olfileid{nml}{syn}{lan}

\olsection{The Language of Basic Modal Logic}

\begin{defn}
The basic language of modal logic contains
\begin{enumerate}
  \tagitem{prvFalse}{The propositional constant for !!{falsity}~$\lfalse$.}{}
  \tagitem{prvTrue}{The propositional constant for !!{truth}~$\ltrue$.}{}
  \item A !!{denumerable}s set of !!{propositional variable}s: $\Obj
    p_0$, $\Obj p_1$, $\Obj p_2$, \dots
  \item The propositional connectives: \startycommalist
  \iftag{prvNot}{\ycomma $\lnot$ (negation)}{}%
  \iftag{prvAnd}{\ycomma $\land$ (conjunction)}{}%
  \iftag{prvOr}{\ycomma $\lor$ (disjunction)}{}%
  \iftag{prvIf}{\ycomma $\lif$ (!!{conditional})}{}%
  \iftag{prvIff}{\ycomma $\liff$ (!!{biconditional})}{}.
  \tagitem{prvBox}{The modal operator $\Box$.}{}
  \tagitem{prvDiamond}{The modal operator $\Diamond$.}{}
\end{enumerate}
\end{defn}

\begin{defn}
\emph{!!^{formula}s} of the basic modal language are inductively
  defined as follows:
\begin{enumerate}
\tagitem{prvFalse}{$\lfalse$ is an atomic !!{formula}.}{}

\tagitem{prvTrue}{$\ltrue$ is an atomic !!{formula}.}{}

\item Every !!{propositional variable} $\Obj p_i$ is an (atomic) !!{formula}.

\tagitem{prvNot}{If $!A$ is !!a{formula}, then $\lnot !A$ is
  !!a{formula}.}{}

\tagitem{prvAnd}{If $!A$ and $!B$ are !!{formula}s, then $(!A \land
  !B)$ is !!a{formula}.}{}

\tagitem{prvOr}{If $!A$ and $!B$ are !!{formula}s, then $(!A \lor !B)$
  is !!a{formula}.}{}

\tagitem{prvIf}{If $!A$ and $!B$ are !!{formula}s, then $(!A \lif !B)$
  is !!a{formula}.}{}

\tagitem{prvIff}{If $!A$ and $!B$ are !!{formula}s, then $(!A \liff !B)$
  is !!a{formula}.}{}

\tagitem{prvBox}{If $!A$ is !!a{formula}, then $\Box !A$ is
  !!a{formula}.}{}

\tagitem{prvDiamond}{If $!A$ is !!a{formula}, then $\Diamond !A$ is
  !!a{formula}.}{}

\tagitem{limitClause}{Nothing else is !!a{formula}.}{}
\end{enumerate}
\end{defn}

\begin{tagblock}{defNot,defOr,defAnd,defIf,defIff,defTrue,defFalse,defBox,defDiamond}
\begin{defn}
Formulas constructed using the defined operators are to be understood
as follows:

\begin{tagenumerate}{defTrue,defFalse,defNot,defOr,defAnd,defIf,defIff,defBox,defDiamond}

\tagitem{defTrue}{$\ltrue$ abbreviates
  \iftag{prvFalse}{$\lnot\lfalse$}{$(!A \lor \lnot !A)$ for some
    fixed atomic !!{formula}~$!A$}.}{}

\tagitem{defFalse}{$\lfalse$ abbreviates
  \iftag{prvTrue}{$\lnot\ltrue$}{$(!A \land \lnot !A)$ for some
    fixed atomic !!{formula}~$!A$}.}{}

\tagitem{defNot}{$\lnot !A$ abbreviates $!A \lif \lfalse$.}{}

\tagitem{defOr}{$!A \lor !B$ abbreviates
  \iftag{prvAnd}{$\lnot(\lnot !A \land \lnot !B)$}{$\lnot !A \lif
    !B$}.}{}

\tagitem{defAnd}{$!A \land !B$ abbreviates
  \iftag{prvOr}{$\lnot(\lnot !A \lor \lnot !B)$}{$\lnot (!A \lif
    \lnot !B)$}.}{}

\tagitem{defIf}{$!A \lif !B$ abbreviates
  \iftag{prvOr}{$\lnot !A \lor !B)$}{$\lnot (!A \land \lnot !B)$}.}{}

\tagitem{defIff}{$!A \liff !B$ abbreviates $(!A \lif !B) \land (!B
  \lif !A)$.}{}

\tagitem{defBox}{$\Box !A$ abbreviates $\lnot\Diamond\lnot !A$. }{}

\tagitem{defDiamond}{$\Diamond !A$ abbreviates $\lnot\Box\lnot !A$.}{}
\end{tagenumerate}
\end{defn}
\end{tagblock}

If !!a{formula}~$!A$ does not contain $\Box$ or $\Diamond$, we say it
is \emph{modal-free}.


% ==== OLP file: content/normal-modal-logic/syntax-and-semantics/modal-validity.tex ====
% Part: normal-modal-logic
% Chapter: syntax-and-semantics
% Section: modal-validity



\olfileid{nml}{syn}{val}

\olsection{Validity}

\begin{explain}
  !!^{formula}s that are true in all models, i.e., true at every world
  in every model, are particularly interesting. They represent those
  modal propositions which are true regardless of how $\Box$ and
  $\Diamond$ are interpreted, as long as the interpretation is
  ``normal'' in the sense that it is generated by some accessibility
  relation on possible worlds. We call such !!{formula}s
  \emph{valid}. For instance, $\Box(p \land q) \lif \Box p$ is
  valid. Some !!{formula}s one might expect to be valid on the basis
  of the alethic interpretation of $\Box$, such as $\Box p \lif p$,
  are not valid, however.  Part of the interest of relational models
  is that different interpretations of $\Box$ and $\Diamond$ can be
  captured by different kinds of accessibility relations. This
  suggests that we should define validity not just relative to
  \emph{all} models, but relative to all models \emph{of a certain
    kind}. It will turn out, e.g., that $\Box p \lif p$ is true in all
  models where every world is accessible from itself, i.e., $R$ is
  reflexive. Defining validity relative to classes of models enables
  us to formulate this succinctly: $\Box p \lif p$ is valid in the
  class of reflexive models.
\end{explain}

\begin{defn}
  !!^a{formula}~$!A$ is \emph{valid} in a class $\mClass{C}$ of
  models if it is true in every model in~$\mClass{C}$ (i.e., true at
  every world in every model in~$\mClass{C}$). If $!A$ is valid
  in~$\mClass{C}$, we write $\mClass{C} \Entails !A$, and we
  write $\Entails !A$ if $!A$ is valid in the class of
  \emph{all} models.
\end{defn}

\begin{prop}\ollabel{prop:subset-class}
  If $!A$ is valid in $\mClass{C}$ it is also valid in each class
  $\mClass{C}' \subseteq \mClass{C}$.
\end{prop}

\begin{prop}\ollabel{prop:Nec-rule}
  If $!A$ is valid, then so is $\Box!A$. 
\end{prop}

\begin{proof}
  Assume $\Entails !A$. To show $\Entails \Box!A$ let $\mModel{M} =
  \tuple{W, R, V}$ be a model and $w \in W$. If $Rww'$ then
  $\mSat{M}{!A}[w']$, since $!A$ is valid, and so also
  $\mSat{M}{\Box!A}[w]$. Since $\mModel{M}$ and $w$ were
  arbitrary, $\Entails \Box!A$.
\end{proof}

\begin{prob}
  Show that the following are valid:
  \begin{enumerate}
  \item $\Box p \lif \Box (q \lif p)$;
  \item $\Box \lnot \lfalse$;
  \item $\Box p \lif (\Box q \lif \Box p)$.
  \end{enumerate}
\end{prob}

\begin{prob}
  Show that $!A \lif \Box!A$ is valid in the class $\mClass{C}$ of
  models $\mModel{M} = \tuple{W, R, V}$ where $W = \{w\}$. Similarly,
  show that $!B \lif \Box !A$ and $\Diamond !A \lif !B$ are valid in
  the class of models $\mModel{M} = \tuple{W, R, V}$ where
  $R = \emptyset$.
\end{prob}


% ==== OLP file: content/normal-modal-logic/syntax-and-semantics/normal-modal-logics.tex ====
% Part: normal-modal-logic
% Chapter: syntax-and-semantics
% Section: normal-modal-logics



\olfileid{nml}{syn}{nor}

\olsection{Normal Modal Logics}

A modal logic can be seen as any suitably well-behaved set of
!!{formula}s.  There are some modal logics which are of particular
interest, both because they have interesting interpretations and
applications, but mainly because they are theoretically well
understood.  Those are the so-called normal modal logics.

\begin{defn}\ollabel{defn:modal-logic}
A \emph{modal logic}~$\Log L$ is a set of !!{formula}s of the basic
modal language which 
\begin{enumerate}
\item contains all propositional tautologies, 
\item is closed under uniform substitution, and 
\item is closed under modus ponens, i.e., whenever $!A$ and $(!A \lif
  !B) \in \Log L$ then also $!B \in \Log L$.
\end{enumerate}
\end{defn}

\begin{ex}
Some admittedly not very interesting sets of !!{formula}s which
satisfy this definition and therefore count as modal logics are:
\begin{enumerate}
\item the set of all tautologies
\item the set of all !!{formula}s (the inconsistent modal logic)
\end{enumerate}
\end{ex}

\begin{defn}\label{defn:normal-modal-logic}
A modal logic~$\Log L$ is \emph{normal} iff 
\begin{enumerate}
\item it contains both
\begin{align}
\tag{K} \Box(p \lif q) & \lif (\Box p \lif \Box q) \\
\tag{Dual} \Diamond p & \liff \lnot \Box \lnot p
\end{align}
\item and it is closed under necessitation, i.e., whenever $!A \in
  \Log L$, then also $\Box !A \in \Log L$.
\end{enumerate}
\end{defn}

We will encounter two main ways of defining modal logics.  One is
proof-theoretically, as the set of !!{formula}s that can be derived in
certain !!{derivation} systems. The other is model-theoretically, as
sets of !!{formula}s which are valid in certain classes of
models. This mirrors the situation in ordinary propositional logic (or
first-order logic, for that matter).  Here we can also single out a
set of !!{formula}s in two different ways, e.g., as the set of
tautologies which are true in all truth-value assignments and as the
set of all !!{derivable} !!{formula}s in some !!{derivation} system.
For normal modal logics, this two-fold way of specifying modal logics
is possible using relational models and axiomatic (Hilbert-type) proof
systems.

% ==== OLP file: content/normal-modal-logic/syntax-and-semantics/relational-models.tex ====
% Part: normal-modal-logic
% Chapter: syntax-and-semantics
% Section: relational-models



\olfileid{nml}{syn}{rel}

\olsection{Relational Models}

The basic concept of semantics for normal modal logics is that of a
\emph{relational model}. It consists of a set of worlds, which are
related by a binary ``accessibility relation,'' together with an
assignment which determines which !!{propositional variable}s count as
``true'' at which worlds.

\begin{defn}
  A \emph{model} for the basic modal language is a triple $\mModel{M}
  = \tuple{W, R, V}$, where
  \begin{enumerate}
  \item $W$ is a nonempty set of ``worlds,''
  \item $R$ is a binary accessibility relation on~$W$, and
  \item $V$ is a function assigning to each !!{propositional
    variable}~$p$ a set $V(p)$ of possible worlds.
  \end{enumerate}
  When $Rww'$ holds, we say that $w'$ is \emph{accessible
    from}~$w$. When $w \in V(p)$ we say $p$ is \emph{true at}~$w$.
\end{defn}

The great advantage of relational semantics is that models can be
represented by means of simple diagrams, such as the one in
\olref{fig:simple}. Worlds are represented by nodes, and world~$w'$ is
accessible from~$w$ precisely when there is an arrow from~$w$
to~$w'$. Moreover, we label a node (world) by~$\mTrue{p}$ when $w \in V(p)$,
and otherwise by~\mFalse{p}. \olref{fig:simple} represents the model
with $W = \{w_1, w_2, w_3\}$, $R = \{\tuple{w_1, w_2},\tuple{w_1,
  w_3}\}$, $V(p) = \{w_1, w_2\}$, and $V(q) = \{w_2\}$.

\begin{figure}
  \begin{center}
    \begin{tikzpicture}[modal]
      \node[world] (w1) [label={[align=right]right:\mTrue{p}\\ \mFalse{q}}]{$w_1$}; 
      \node[world] (w2) [label={[align=right]right:\mTrue{p}\\ \mTrue{q}},
        above right=of w1]{$w_2$}; 
      \node[world] (w3) [label={[align=right]right:\mFalse{p}\\ \mFalse{q}},
        below right=of w1] {$w_3$};
      \draw[->] (w1) to (w2);
      \draw[->] (w1) to (w3);
    \end{tikzpicture}
  \end{center}
  \caption{A simple model.}
  \ollabel{fig:simple}
\end{figure}


% ==== OLP file: content/normal-modal-logic/syntax-and-semantics/schemas.tex ====
% Part: normal-modal-logic
% Chapter: syntax-and-semantics
% Section: schemas



\olfileid{nml}{syn}{sch}

\olsection{Schemas and Validity}

\begin{defn}
  A \emph{schema} is a set of !!{formula}s comprising all
  and only the substitution instances of some modal !!{formula}~$!C$, i.e.,
  \[
  \Setabs{!B}{\lexists[!D_1], \dots, \lexists[!D_n] \left(!B =
  \SSubst{!C}{\subst{!D_1}{p_1}, \dots, \subst{!D_n}{p_n}} \right) }.
  \]
  The !!{formula}~$!C$ is called the \emph{characteristic} !!{formula} of
  the schema, and it is unique up to a renaming of the !!{propositional
  variable}s. !!^a{formula}~$!A$ is an \emph{instance} of a schema if
  it is a member of the set.
\end{defn}

It is convenient to denote a schema by the meta-linguistic expression
obtained by substituting `$!A$', `$!B$',~\dots, for the atomic
components of~$!C$. So, for instance, the following denote schemas:
`$!A$', `$!A \lif \Box!A$', `$!A \lif (!B \lif !A)$'. They correspond
to the characteristic !!{formula}s $p$, $p \lif \Box p$, $p \lif (q
\lif p)$. The schema `$!A$' denotes the set of \emph{all}
!!{formula}s.

\begin{defn}
  A schema is \emph{true} in a model if and only if all of its instances
  are; and a schema is \emph{valid} if and only if it is true in every
  model.
\end{defn}

\begin{prop}\ollabel{prop:Kvalid}
  The following schema \Ax{K} is valid
  \begin{equation*}
  \Box(!A \lif !B) \lif (\Box !A \lif \Box !B). \tag{\Ax{K}}
  \end{equation*}
\end{prop}

\begin{proof}
  We need to show that all instances of the schema are true at every
  world in every model. So let $\mModel{M} = \tuple{W,R,V}$ and $w \in
  W$ be arbitrary. To show that a conditional is true at a world we
  assume the antecedent is true to show that the consequent is true as
  well. In this case, let $\mSat{M}{\Box(!A \lif !B)}[w]$ and
  $\mSat{M}{\Box !A}[w]$. We need to show $\mSat{M}{\Box !B}[w]$. So let
  $w'$ be arbitrary such that $Rww'$. Then by the first assumption
  $\mSat{M}{!A \lif !B}[w']$ and by the second assumption
  $\mSat{M}{!A}[w']$. It follows that $\mSat{M}{!B}[w']$. Since $w'$
  was arbitrary, $\mSat{M}{\Box !B}[w]$.
\end{proof}

\begin{prop}\ollabel{prop:Dual-valid}
  The following schema \Dual{} is valid
  \begin{equation*}
  \Diamond !A \liff \lnot\Box\lnot !A. \tag{\Dual}
  \end{equation*}
\end{prop}

\begin{proof}
  Exercise.
\end{proof}

\begin{prob}
  Prove \olref[nml][syn][sch]{prop:Dual-valid}.
\end{prob}

\begin{prop}\ollabel{prop:soundMP}
  If $!A$ and $!A \lif !B$ are true at a world in a model
  then so is $!B$. Hence, the valid !!{formula}s are closed under
  modus ponens.
\end{prop}

\begin{prop}\ollabel{prop:valid-instances}
  !!^a{formula}~$!A$ is valid iff all its substitution instances
  are. In other words, a schema is valid iff its characteristic
  !!{formula} is.
\end{prop}

\begin{proof}
  The ``if'' direction is obvious, since $!A$ is a substitution
  instance of itself.

  To prove the ``only if'' direction, we show the following: Suppose
  $\mModel{M} = \tuple{W, R, V}$ is a modal model, and $!B \ident
  \SSubst{!A}{\subst{!D_1}{p_1},\dots,\subst{!D_n}{p_n}}$ is a
  substitution instance of~$!A$. Define $\mModel{M'} = \tuple{W, R,
    V'}$ by $V'(p_i) = \Setabs{w}{\mSat{M}{!D_i}[w]}$.  Then
  $\mSat{M}{!B}[w]$ iff $\mSat{M'}{!A}[w]$, for any~$w \in W$. (We
  leave the proof as an exercise.) Now suppose that $!A$ was valid,
  but some substitution instance $!B$ of~$!A$ was not valid. Then for
  some $\mModel{M} = \tuple{W, R, V}$ and some $w \in W$,
  $\mSat/{M}{!B}[w]$. But then $\mSat/{M'}{!A}[w]$ by the claim, and
  $!A$ is not valid, a contradiction.
\end{proof}

\begin{prob}
  Prove the claim in the ``only if'' part of the proof of
  \olref[nml][syn][sch]{prop:valid-instances}. (Hint: use induction
  on~$!A$.)
\end{prob}

Note, however, that it is not true that a schema is true in a model
iff its characteristic formula is. Of course, the ``only if''
direction holds: if every instance of $!A$ is true in~$\mModel{M}$,
$!A$~itself is true in~$\mModel{M}$. But it may happen that $!A$ is
true in~$\mModel{M}$ but some instance of~$!A$ is false at some world
in~$\mModel{M}$. For a very simple counterexample consider~$p$ in a
model with only one world~$w$ and $V(p) = \{w\}$, so that $p$ is true
at~$w$. But $\lfalse$ is an instance of~$p$, and not true at~$w$.

\begin{prob}
  Show that none of the following !!{formula}s are valid:
  \begin{enumerate}
    \item[\Ax{D}:] \quad $\Box p \lif \Diamond p$;
    \item[\Ax{T}:] \quad $\Box p \lif p$;
    \item[\Ax{B}:] \quad $p \lif \Box\Diamond p$;
    \item[\Ax{4}:] \quad $\Box p \lif \Box \Box p$;
    \item[\Ax{5}:] \quad $\Diamond p \lif \Box \Diamond
      p$.
  \end{enumerate}
\end{prob}

\begin{table}[t]
    \centering
    \begin{tabular}{| l || l |}
      \hline
      {\emph{Valid Schemas}} & {\emph{Invalid Schemas}} \\
      \hline\hline
      $\Box(!A \lif !B) \lif (\Diamond !A \lif \Diamond !B)$
      & $\Box (!A \lor !B) \lif (\Box !A \lor \Box !B)$ \\
      $\Diamond (!A \lif !B) \lif (\Box !A \lif \Diamond
      !B)$
      & $(\Diamond !A \land \Diamond !B) \lif \Diamond (!A
      \land !B)$\\
      $\Box (!A \land !B) \liff (\Box !A \land \Box !B)$
      & $!A \lif \Box !A$ \\
      $\Box !A \lif \Box (!B \lif !A)$
      & $\Box \Diamond !A \lif !B$ \\
      $\lnot \Diamond !A \lif \Box (!A \lif !B)$
      & $\Box \Box !A \lif \Box !A$ \\
      $\Diamond (!A \lor !B) \liff (\Diamond !A \lor
      \Diamond !B)$
      & $\Box \Diamond !A \lif \Diamond \Box !A$. \\
      \hline
    \end{tabular}
    \caption{Valid and (or?) invalid schemas.}
    \ollabel{tab:valid-invalidSchemas}
\end{table}

\begin{prob}%\ollabel{ex:in/validSchemas}
  Prove that the schemas in the first column of
  \olref[nml][syn][sch]{tab:valid-invalidSchemas} are valid and those in
  the second column are not valid.
\end{prob}

\begin{prob}
  Decide whether the following schemas are valid or invalid:
  \begin{enumerate}
  \item $(\Diamond !A \lif \Box !B) \lif (\Box !A \lif \Box
    !B)$;
  \item $\Diamond(!A \lif !B) \lor \Box(!B \lif !A)$.
  \end{enumerate}
\end{prob}

\begin{prob}
  For each of the following schemas find a model $\mModel{M}$ such
  that every instance of the !!{formula} is true in $\mModel{M}$:
  \begin{enumerate}
  \item $p \lif \Diamond\Diamond p$;
  \item $\Diamond p \lif \Box p$.
  \end{enumerate}
\end{prob}


% ==== OLP file: content/normal-modal-logic/syntax-and-semantics/substitution.tex ====
% Part: normal-modal-logic
% Chapter: syntax-and-semantics
% Section: substitution



\olfileid{nml}{syn}{sub}

\olsection{Simultaneous Substitution}

An \emph{instance} of !!a{formula}~$!A$ is the result of replacing all
occurrences of !!a{propositional variable} in~$!A$ by some other
!!{formula}. We will refer to instances of !!{formula}s often, both
when discussing validity and when discussing !!{derivability}. It
therefore is useful to define the notion precisely.

\begin{defn}\ollabel{def:subst-inst}
  Where $!A$ is a modal !!{formula} all of whose !!{propositional
    variable}s are among $p_1$, \dots, $p_n$, and $!D_1$, \dots,
  $!D_n$ are also modal !!{formula}s, we define
  $\SSubst{!A}{\subst{!D_1}{p_1}, \dots, \subst{!D_n}{p_n}}$ as the
  result of simultaneously substituting each $!D_i$ for $p_i$
  in~$!A$. Formally, this is a definition by induction on~$!A$:
  \begin{enumerate}
    \tagitem{prvFalse}{\indcase{!A}{\lfalse}{%
        $\SSubst{\indfrm}{\subst{!D_1}{p_1}, \dots,
          \subst{!D_n}{p_n}}$ is $\lfalse$.}}{}
    \tagitem{prvTrue}{\indcase{!A}{\ltrue}{%
        $\SSubst{\indfrm}{\subst{!D_1}{p_1}, \dots,
          \subst{!D_n}{p_n}}$ is $\ltrue$.}}{}
    \item \indcase{!A}{q}{$\SSubst{\indfrm}{\subst{!D_1}{p_1}, \dots,
        \subst{!D_n}{p_n}}$ is $q$, provided $q \not\ident p_i$ for $i
      = 1$, \dots,~$n$.}
    \item \indcase{!A}{p_i}{$\SSubst{\indfrm}{\subst{!D_1}{p_1},
        \dots, \subst{!D_n}{p_n}}$ is $!D_i$.}
    \tagitem{prvNot}{\indcase{!A}{\lnot !B}{%
        $\SSubst{\indfrm}{\subst{!D_1}{p_1}, \dots, \subst{!D_n}{p_n}}$ is
        $\lnot \SSubst{!B}{\subst{!D_1}{p_1}, \dots, \subst{!D_n}{p_n}}$.}}{}
    \tagitem{prvAnd}{\indcase{!A}{(!B \land
        !C)}{$\SSubst{\indfrm}{\subst{!D_1}{p_1}, \dots,
          \subst{!D_n}{p_n}}$ is
        \[(\SSubst{!B}{\subst{!D_1}{p_1}, \dots, \subst{!D_n}{p_n}} \land
          \SSubst{!C}{\subst{!D_1}{p_1}, \dots, \subst{!D_n}{p_n}}).\]}}{}
    \tagitem{prvOr}{\indcase{!A}{(!B \lor
          !C)}{$\SSubst{\indfrm}{\subst{!D_1}{p_1}, \dots,
          \subst{!D_n}{p_n}}$ is
        \[(\SSubst{!B}{\subst{!D_1}{p_1}, \dots, \subst{!D_n}{p_n}} \lor
          \SSubst{!C}{\subst{!D_1}{p_1}, \dots, \subst{!D_n}{p_n}}).\]}}{}
    \tagitem{prvIf}{\indcase{!A}{(!B \lif
          !C)}{$\SSubst{\indfrm}{\subst{!D_1}{p_1}, \dots,
          \subst{!D_n}{p_n}}$ is
        \[(\SSubst{!B}{\subst{!D_1}{p_1}, \dots, \subst{!D_n}{p_n}} \lif
          \SSubst{!C}{\subst{!D_1}{p_1}, \dots, \subst{!D_n}{p_n}}).\]}}{}
    \tagitem{prvIf}{\indcase{!A}{(!B \liff
          !C)}{$\SSubst{\indfrm}{\subst{!D_1}{p_1}, \dots,
          \subst{!D_n}{p_n}}$ is
        \[(\SSubst{!B}{\subst{!D_1}{p_1}, \dots, \subst{!D_n}{p_n}} \liff
          \SSubst{!C}{\subst{!D_1}{p_1}, \dots, \subst{!D_n}{p_n}}).\]}}{}
    \item \indcase{!A}{\Box !B}{%
        $\SSubst{\indfrm}{\subst{!D_1}{p_1}, \dots, \subst{!D_n}{p_n}}$ is
        $\Box
        \SSubst{!B}{\subst{!D_1}{p_1}, \dots, \subst{!D_n}{p_n}}.$}
    \tagitem{prvDiamond}{\indcase{!A}{\Diamond !B}{%
        $\SSubst{\indfrm}{\subst{!D_1}{p_1}, \dots, \subst{!D_n}{p_n}}$ is
        $\Diamond
        \SSubst{!B}{\subst{!D_1}{p_1}, \dots, \subst{!D_n}{p_n}}.$}}{}
  \end{enumerate}
  The !!{formula} $\SSubst{!A}{\subst{!D_1}{p_1}, \dots,
    \subst{!D_n}{p_n}}$ is called a \emph{substitution instance}
  of~$!A$.
\end{defn}

\begin{ex}
  Suppose $!A$ is $p_1 \lif \Box(p_1 \land p_2)$, $!D_1$ is
  $\Diamond(p_2 \lif p_3)$ and $!D_2$ is $\lnot\Box p_1$. Then
  $\SSubst{!A}{\subst{!D_1}{p_1}, \subst{!D_2}{p_2}}$ is
  \begin{align*}
    \Diamond(p_2 \lif p_3) &
    \lif \Box(\Diamond(p_2 \lif p_3) \land \lnot\Box p_1)
    \intertext{while $\SSubst{!A}{\subst{!D_2}{p_1}, \subst{!D_1}{p_2}}$ is}
    \lnot\Box p_1 &
    \lif \Box(\lnot\Box p_1 \land \Diamond(p_2 \lif p_3))
    \intertext{Note that simultaneous substitution is in general not
      the same as iterated substitution, e.g., compare
      $\SSubst{!A}{\subst{!D_1}{p_1}, \subst{!D_2}{p_2}}$ above with
      $\Subst{(\Subst{!A}{!D_1}{p_1})}{!D_2}{p_2}$, which is:}
    \Diamond(p_2 \lif p_3) & \Subst{\lif \Box(\Diamond(p_2 \lif p_3) \land p_2)}{\lnot\Box p_1}{p_2}, \text{ i.e.,}\\
    \Diamond(\lnot\Box p_1 \lif p_3) &
    \lif \Box(\Diamond(\lnot\Box p_1
    \lif p_3) \land \lnot\Box p_1)\\
    \intertext{and with $\Subst{(\Subst{!A}{!D_2}{p_2})}{!D_1}{p_1}$:}
    p_1 & \lif \Subst{\Box(p_1 \land \lnot\Box p_1)}{\Diamond(p_2 \lif p_3)}{p_1}, \text{ i.e.,}\\
    \Diamond(p_2 \lif p_3) &
    \lif \Box(\Diamond(p_2
    \lif p_3) \land \lnot\Box\Diamond(p_2 \lif p_3)).
  \end{align*}
\end{ex}


% ==== OLP file: content/normal-modal-logic/syntax-and-semantics/syntax-and-semantics.tex ====
% Part: normal-modal-logic
% Chapter: syntax-and-semantics



\olchapter{nml}{syn}{Syntax and Semantics}

\olimport{introduction}
\olimport{language-modal-logic}
\olimport{substitution}
\olimport{relational-models}
\olimport{truth-at-w}
\olimport{truth-in-model}
\olimport{modal-validity}
\olimport{tautological-instances}
\olimport{schemas}
\olimport{entailment}

\OLEndChapterHook



% ==== OLP file: content/normal-modal-logic/syntax-and-semantics/tautological-instances.tex ====
% Part: normal-modal-logic
% Chapter: syntax-and-semantics
% Section: tautological-instances



\olfileid{nml}{syn}{tau}

\olsection{Tautological Instances}

\begin{explain}
A modal-free formula is a tautology if it is true under every
truth-value assignment. Clearly, every tautology is true at every
world in every model. But for !!{formula}s involving $\Box$
and~$\Diamond$, the notion of tautology is not defined. Is it the
case, e.g., that $\Box p \lor \lnot \Box p$---an instance of the
principle of excluded middle---is valid? The notion of a
\emph{tautological instance} helps: !!a{formula} that is a
substitution instance of a (non-modal) tautology. It is not
surprising, but still requires proof, that every tautological instance
is valid.
\end{explain}

\begin{defn}
  A modal !!{formula}~$!B$ is a \emph{tautological instance} if and
  only if there is a modal-free tautology~$!A$ with !!{propositional
    variable}s $p_1$, \dots,~$p_n$ and !!{formula}s $!D_1$,
  \dots,~$!D_n$ such that $!B \ident \SSubst{!A}{\subst{!D_1}{p_1},
    \dots, \subst{!D_n}{p_n}}$.
\end{defn}

\begin{lem}\ollabel{lem:valid-taut}
  Suppose $!A$ is a modal-free !!{formula} whose !!{propositional
  variable}s are $p_1$, \dots, $p_n$, and let $!D_1$, \dots,
  $!D_n$ be modal !!{formula}s. Then for any assignment $\pAssign{v}$,
  any model $\mModel{M} = \tuple{W, R, V}$, and any $w \in W$ such
  that $\pAssign{v}(p_i) = \True$ if and only if $\mSat{M}{!D_i}[w]$ we have
  that $\pSat{v}{!A}$ if and only if
  $\mSat{M}{\SSubst{!A}{\subst{!D_1}{p_1}, \dots,
      \subst{!D_n}{p_n}}}[w]$.
\end{lem}

\begin{proof}
  By induction on $!A$.
  \begin{enumerate}
  \tagitem{prvFalse}{\indcase{!A}{\lfalse}{Both $\pSat/{v}{\lfalse}$
      and $\mSat/{M}{\lfalse}[w]$}.}{}
  \tagitem{prvTrue}{\indcase{!A}{\ltrue}{Both $\pSat{v}{\ltrue}$
      and $\mSat{M}{\ltrue}[w]$}.}{}
  \item \indcase{!A}{p_i}{%
    \begin{align*}
      \pSat{v}{p_i} \Leftrightarrow {} & \pAssign{v}(p_i) = \True \\
      & \qquad \text{by definition of $\pSat{v}{p_i}$}\\
      \Leftrightarrow {} & \mSat{M}{!D_i}[w] \\
      &\qquad \text{by assumption}\\
      \Leftrightarrow {} & \mSat{M}{\SSubst{p_i}{\subst{!D_1}{p_1}, \dots,
          \subst{!D_n}{p_n}}}[w]\\
      &\qquad \text{since $\SSubst{p_i}{\subst{!D_1}{p_1}, \dots,
          \subst{!D_n}{p_n}} \ident !D_i$}.
      \end{align*}}
  \tagitem{prvFalse}{\indcase{!A}{\lnot !B}{%
      \begin{align*}
        \pSat{v}{\lnot !B} \Leftrightarrow {} & \pSat/{v}{!B}\\
        &\qquad \text{by definition of $\pSat{v}{}$};\\
        \Leftrightarrow {} & \mSat/{M}{\SSubst{!B}{\subst{!D_1}{p_1}, \dots,
          \subst{!D_n}{p_n}}}[w]\\
        &\qquad \text{by induction hypothesis}\\
        \Leftrightarrow {} &
        \mSat{M}{\SSubst{\lnot !B}{\subst{!D_1}{p_1}, \dots,
          \subst{!D_n}{p_n}}}[w]\\
        &\qquad \text{by definition of $\pSat{v}{}$}.
      \end{align*}}}{}
  \tagitem{prvAnd}{\indcase{!A}{(!B \land !C)}{%
      \begin{align*}
        \pSat{v}{!B \land !C} \Leftrightarrow {} &
        \pSat{v}{!B} \text{ and } \pSat{v}{!C}\\
        &\qquad \text{by definition of $\pSat{v}{}$}\\
        \Leftrightarrow {} &
        \mSat{M}{\SSubst{!B}{\subst{!D_1}{p_1}, \dots,
            \subst{!D_n}{p_n}}}[w] \text{ and } \\
        & \mSat{M}{\SSubst{!C}{\subst{!D_1}{p_1}, \dots,
          \subst{!D_n}{p_n}}}[w]\\
        &\qquad \text{by induction hypothesis}\\
        \Leftrightarrow{} &
        \mSat{M}{\SSubst{(!B \land !C)}{\subst{!D_1}{p_1}, \dots,
          \subst{!D_n}{p_n}}}[w]\\
        &\qquad \text{by definition of $\mSat{M}{}[w]$}.
      \end{align*}}}{}
  \tagitem{prvOr}{\indcase{!A}{(!B \lor !C)}{%
      \begin{align*}
        \pSat{v}{!B \lor !C} \Leftrightarrow{} &
        \pSat{v}{!B} \text{ or } \pSat{v}{!C}\\
        &\qquad \text{by definition of $\pSat{v}{}$};\\
        \Leftrightarrow{} & \mSat{M}{\SSubst{!B}{\subst{!D_1}{p_1}, \dots,
            \subst{!D_n}{p_n}}}[w] \text{ or }\\
        & \mSat{M}{\SSubst{!C}{\subst{!D_1}{p_1}, \dots,
          \subst{!D_n}{p_n}}}[w]\\
        &\qquad \text{by induction hypothesis}\\
        \Leftrightarrow{} &
        \mSat{M}{\SSubst{(!B \lor !C)}{\subst{!D_1}{p_1}, \dots,
          \subst{!D_n}{p_n}}}[w]\\
        &\qquad \text{by definition of $\mSat{M}{}[w]$}.
      \end{align*}}}{}
  \tagitem{prvIf}{\indcase{!A}{(!B \lif !C)}{%
      \begin{align*}
        \pSat{v}{!B \lif !C} \Leftrightarrow{} &
        \pSat/{v}{!B} \text{ or } \pSat{v}{!C}\\
        &\qquad \text{by definition of $\pSat{v}{}$}\\
        \Leftrightarrow{} &
        \mSat/{M}{\SSubst{!B}{\subst{!D_1}{p_1}, \dots,
            \subst{!D_n}{p_n}}}[w] \text{ or }\\
        & \mSat{M}{\SSubst{!C}{\subst{!D_1}{p_1}, \dots,
          \subst{!D_n}{p_n}}}[w]\\
        &\qquad \text{by induction hypothesis}\\
        \Leftrightarrow{} &
        \mSat{M}{\SSubst{(!B \lif !C)}{\subst{!D_1}{p_1}, \dots,
          \subst{!D_n}{p_n}}}[w]\\
        &\qquad \text{by definition of $\mSat{M}{}[w]$}.
      \end{align*}}}{}
  \tagitem{prvIff}{\indcase{!A}{(!B \liff !C)}{%
      \begin{align*}
        \pSat{v}{!B \lif !C} \Leftrightarrow {} &
        \text{either } \pSat{v}{!B} \text{ and } \pSat{v}{!C}\\
        & \text{or } \pSat/{v}{!B} \text{ and } \pSat/{v}{!C}\\
        &\qquad \text{by definition of $\pSat{v}{}$}\\
        \Leftrightarrow {} &
        \text{either } \mSat{M}{\SSubst{!B}{\subst{!D_1}{p_1}, \dots,
            \subst{!D_n}{p_n}}}[w] \text{ and }\\
        & \qquad \mSat{M}{\SSubst{!C}{\subst{!D_1}{p_1}, \dots,
          \subst{!D_n}{p_n}}}[w]\\ 
        &  \text{ or } \mSat/{M}{\SSubst{!B}{\subst{!D_1}{p_1}, \dots,
            \subst{!D_n}{p_n}}}[w] \text{ and }\\
        & \qquad \mSat/{M}{\SSubst{!C}{\subst{!D_1}{p_1}, \dots,
          \subst{!D_n}{p_n}}}[w]\\
        &\qquad\text{by induction hypothesis}\\
        \Leftrightarrow {} &
        \mSat{M}{\SSubst{(!B \liff !C)}{\subst{!D_1}{p_1}, \dots,
          \subst{!D_n}{p_n}}}[w]\\
        &\qquad \text{by definition of $\mSat{M}{}[w]$}.
      \end{align*}}}{}
  \end{enumerate}
  \end{proof}

\begin{prop}\ollabel{prop:valid-taut}
  All tautological instances are valid.
\end{prop}

\begin{proof}
  Contrapositively, suppose $!A$ is such that
  $\mSat/{M}{\SSubst{!A}{\subst{!D_1}{p_1}, \dots,
  \subst{!D_n}{p_n}}}[w]$, for some model $\mModel{M}$ and world~$w$.
  Define an assignment $\pAssign{v}$ such that $\pAssign{v}(p_i) =
  \True$ if and only if $\mSat{M}{!D_i}[w]$ (and $\pAssign{v}$ assigns
  arbitrary values to $q \notin \{p_1, \dots, p_n \}$).  Then by
  \olref{lem:valid-taut}, $\pSat/{v}{!A}$, so $!A$ is not a tautology.
\end{proof}


% ==== OLP file: content/normal-modal-logic/syntax-and-semantics/truth-at-w.tex ====
% Part: normal-modal-logic
% Chapter: syntax-and-semantics
% Section: truth-at-w



\olfileid{nml}{syn}{trw}

\olsection{Truth at a World}

Every modal model determines which modal !!{formula}s count as true at
which worlds in it. The relation ``model $\mModel{M}$ makes
!!{formula}~$!A$ true at world~$w$'' is the basic notion of relational
semantics. The relation is defined inductively and coincides with the
usual characterization using truth tables for the non-modal operators.

\begin{defn}\ollabel{defn:mmodels}
  \emph{Truth of !!a{formula}~$!A$ at~$w$} in a~$\mModel M$, in symbols:
  $\mSat{M}{!A}[w]$, is defined inductively as follows:
  \begin{enumerate}
  \tagitem{prvFalse}{\indcase{!A}{\lfalse}{Never $\mSat{M}{\lfalse}[w]$}.}{}
  \tagitem{prvTrue}{\indcase{!A}{\ltrue}{Always $\mSat{M}{\ltrue}[w]$}.}{}
  \item $\mSat{M}{p}[w]$ iff $w \in V(p)$.
  \tagitem{prvNot}{\indcase{!A}{\lnot !B}{$\mSat{M}{\indfrm}[w]$ iff
    $\mSat/{M}{!B}[w]$}.}{}
  \tagitem{prvAnd}{\indcase{!A}{(!B \land !C)}{$\mSat{M}{\indfrm}[w]$ iff
    $\mSat{M}{!B}[w]$ and $\mSat{M}{!C}[w]$}.}{}
  \tagitem{prvOr}{\indcase{!A}{(!B \lor !C)}{$\mSat{M}{\indfrm}[w]$ iff
    $\mSat{M}{!B}[w]$ or $\mSat{M}{!C}[w]$} (or both).}{}
  \tagitem{prvIf}{\indcase{!A}{(!B \lif !C)}{$\mSat{M}{\indfrm}[w]$ iff
    $\mSat/{M}{!B}[w]$ or $\mSat{M}{!C}[w]$}.}{}
  \tagitem{prvIff}{\indcase{!A}{(!B \liff !C)}{$\mSat{M}{\indfrm}[w]$ iff
    either both $\mSat{M}{!B}[w]$ and $\mSat{M}{!C}[w]$ or
    neither $\mSat{M}{!B}[w]$ nor $\mSat{M}{!C}[w]$}.}{}
  \item\ollabel{defn:sub:mmodels-box}
    \indcase{!A}{\Box !B}{$\mSat{M}{\indfrm}[w]$ iff
    $\mSat{M}{!B}[w']$ for all $w' \in W$ with $Rww'$.}
  \item\ollabel{defn:sub:mmodels-diamond}
    \indcase{!A}{\Diamond !B}{$\mSat{M}{\indfrm}[w]$ iff
    $\mSat{M}{!B}[w']$ for at least one $w' \in W$ with $Rww'$.}
  \end{enumerate}
\end{defn}

Note that by clause~\olref{defn:sub:mmodels-box}, !!a{formula}~$\Box
!B$ is true at~$w$ whenever there are no~$w'$ with $Rww'$. In
such a case $\Box !B$ is \emph{vacuously} true at~$w$. Also,
$\Box !B$ may be satisfied at~$w$ even if $!B$ is not. The truth
of~$!B$ at~$w$ does not guarantee the truth of~$\Diamond !B$
at~$w$. This holds, however, if $Rww$, e.g., if $R$ is reflexive.  If
there is no $w'$ such that $Rww'$, then $\mSat/{M}{\Diamond !A}[w]$,
for any~$!A$.

\begin{prob}
  Consider the model of \olref[nml][syn][rel]{fig:simple}. Which of the
  following hold?
    \begin{enumerate}
    \item $\mSat{M}{q}[w_1]$;
    \item $\mSat{M}{\lnot q}[w_3]$;
    \item $\mSat{M}{p \lor q}[w_1]$;
    \item $\mSat{M}{\Box (p \lor q)}[w_1]$;
    \item $\mSat{M}{\Box q}[w_3]$;
    \item $\mSat{M}{\Box \bot}[w_3]$;
    \item $\mSat{M}{\Diamond q}[w_1]$;
    \item $\mSat{M}{\Box q}[w_1]$;
    \item $\mSat{M}{\lnot \Box\Box \lnot q}[w_1]$.
    \end{enumerate}
\end{prob}

% Prove duality of \Box and \Diamond
\begin{tagblock}{prvDiamond}
\begin{prop}\ollabel{prop:dual}
  \begin{enumerate}
  \item $\mSat{M}{\Box !A}[w]$ iff $\mSat{M}{\lnot\Diamond\lnot !A}[w]$.
  \item $\mSat{M}{\Diamond !A}[w]$ iff $\mSat{M}{\lnot\Box\lnot !A}[w]$.
  \end{enumerate}
\end{prop}

\begin{proof}
  \begin{enumerate}
    \item $\mSat{M}{\lnot\Diamond\lnot !A}[w]$ iff $\mSat/{M}{\Diamond\lnot
      !A}[w]$ by definition of $\mSat{M}{}[w]$. $\mSat{M}{\Diamond\lnot
      !A}[w]$ iff for some $w'$ with $Rww'$, $\mSat{M}{\lnot
      !A}[w']$. Hence, $\mSat/{M}{\Diamond\lnot !A}[w]$ iff for all
    $w'$ with $Rww'$, $\mSat/{M}{\lnot !A}[w']$. We also have
    $\mSat/{M}{\lnot !A}[w']$ iff $\mSat{M}{!A}[w']$. Together we have
    $\mSat{M}{\lnot\Diamond\lnot !A}[w]$ iff for all $w'$ with $Rww'$,
    $\mSat{M}{!A}[w']$. Again by definition of $\mSat{M}{}[w]$, that
      is the case iff $\mSat{M}{\Box !A}[w]$.
    \item \iftag{probDiamond}{Exercise.}{$\mSat{M}{\lnot\Box\lnot !A}[w]$ iff
      $\mSat/{M}{\Box\lnot !A}$. $\mSat{M}{\Box\lnot !A}[w]$ iff for all
      $w'$ with $Rww'$, $\mSat{M}{\lnot !A}[w']$. Hence,
      $\mSat/{M}{\Box\lnot !A}[w]$ iff for some $w'$ with $Rww'$,
      $\mSat/{M}{\lnot !A}[w']$. We also have $\mSat/{M}{\lnot
        !A}[w']$ iff $\mSat{M}{!A}[w']$. Together we have
      $\mSat{M}{\lnot\Box\lnot !A}[w]$ iff for some~$w'$ with
      $Rww'$, $\mSat{M}{!A}[w']$. Again by definition of
      $\mSat{M}{}[w]$, that is the case iff $\mSat{M}{\Diamond
        !A}[w]$.}
  \end{enumerate}
\end{proof}

\begin{probtag}{probDiamond}
  Complete the proof of \olref[nml][syn][trw]{prop:dual}.
\end{probtag}
\end{tagblock}

\begin{prob}
  Let $\mModel{M} = \tuple{W, R, V}$ be a model, and suppose $w_1, w_2 \in
  W$ are such that:
  \begin{enumerate}
  \item $w_1 \in V(p)$ if and only if $w_2 \in V(p)$ (for every !!{propositional variable} $p$); and
  \item for all $w \in W$: $Rw_1w$ if and only if $Rw_2w$.
  \end{enumerate}
  Using induction on !!{formula}s, show that for all !!{formula}s $!A$:
  $\mSat{M}{!A}[w_1]$ if and only if $\mSat{M}{!A}[w_2]$.
\end{prob}

\begin{probtag}{prvDiamond}
  Let $\mModel{M} = \tuple{W, R, V}$. Show that
  $\mSat{M}{\lnot\Diamond !A}[w]$ if and only if
  $\mSat{M}{\Box\lnot!A}[w]$.
\end{probtag}


% ==== OLP file: content/normal-modal-logic/syntax-and-semantics/truth-in-model.tex ====
% Section: truth-in-model



\olfileid{nml}{syn}{tru}

\olsection{Truth in a Model}

Sometimes we are interested in which !!{formula}s are true at every world
in a given model. Let's introduce a notation for this.

\begin{defn}
  !!^a{formula}~$!A$ is \emph{true in a model} $M = \tuple{W, R,
    V}$, written $\mSat{M}{!A}$, if and only if $\mSat{M}{!A}[w]$
  for every $w \in W$.
\end{defn}

\begin{prop}\ollabel{prop:truthfacts}
  \begin{enumerate}
  \item If $\mSat{M}{!A}$ then $\mSat/{M}{\lnot !A}$, but \emph{not}
    vice-versa.
  \item If $\mSat{M}{!A \lif !B}$ then $\mSat{M}{!A}$ only if
    $\mSat{M}{!B}$, but \emph{not} vice-versa.
  \end{enumerate}
\end{prop}

\begin{proof}
  \begin{enumerate}
  \item If $\mSat{M}{!A}$ then $!A$ is true at all worlds in $W$, and
    since $W \neq \emptyset$, it can't be that $\mSat{M}{\lnot!A}$, or
    else $!A$ would have to be both true and false at some
    world.

    On the other hand, if $\mSat/{M}{\lnot !A}$ then $!A$ is true at
    some world $w \in W$. It does not follow that $\mSat{M}{!A}[w]$
    for \emph{every}~$w \in W$. For instance, in the model of
    \olref[rel]{fig:simple}, $\mSat/{M}{\lnot p}$, and also $\mSat/{M}{p}$.
  \item Assume $\mSat{M}{!A \lif !B}$ and $\mSat{M}{!A}$; to show
    $\mSat{M}{!B}$ let $w \in W$ be an arbitrary world. Then
    $\mSat{M}{!A \lif !B}[w]$ and $\mSat{M}{!A}[w]$, so
    $\mSat{M}{!B}[w]$, and since $w$ was arbitrary,
    $\mSat{M}{!B}$.

    To show that the converse fails, we need to find a model
    $\mModel{M}$ such that $\mSat{M}{!A}$ only if $\mSat{M}{!B}$, but
    $\mSat/{M}{!A \lif !B}$. Consider again the model of
    \olref[rel]{fig:simple}: $\mSat/{M}{p}$ and hence (vacuously)
    $\mSat{M}{p}$ only if $\mSat{M}{q}$. However, $\mSat/{M}{p \lif
      q}$, as $p$ is true but $q$ false at~$w_1$.
  \end{enumerate}
\end{proof}

\begin{prob}
  Consider the following model $\mModel{M}$ for the language
  comprising $p_1$, $p_2$, $p_3$ as the only !!{propositional variable}s:
  \begin{center}
    \begin{tikzpicture}[modal]
      \node[world] (w1) [label={[align=right]left:\mTrue{p_1}\\\mFalse{p_2}\\\mFalse{p_3}}]
            {$w_1$} ; 
      \node[world] (w2) [label={[align=right]right:\mTrue{p_1}\\\mTrue{p_2}\\\mFalse{p_3}},
        below right=of w1]  {$w_2$}; 
      \node[world] (w3) [label={[align=right]right:\mTrue{p_1}\\\mTrue{p_2}\\\mTrue{p_3}},
        above right=of w2] {$w_3$};
      \draw[reflexive above] (w3) to (w3); 
      \draw[->] (w1) to (w2);
      \draw[->] (w2) to (w3);
      \draw[->] (w1) to (w3); 
    \end{tikzpicture}
  \end{center}
  Are the following !!{formula}s and schemas true in the model $\mModel{M}$,
  i.e., true at every world in $\mModel{M}$? Explain.
  \begin{enumerate}
  \item $p\lif \Diamond p$ (for $p$ atomic);
  \item $!A\lif \Diamond !A$ (for $!A$ arbitrary);
  \item $\Box p \lif p$ (for $p$ atomic);
  \item $\lnot p \lif \Diamond \Box p$ (for $p$ atomic);
  \item $\Diamond \Box !A$ (for $!A$ arbitrary);
  \item $\Box \Diamond p$ (for $p$ atomic). 
  \end{enumerate}
\end{prob}


